% Vietnamese translation for OI Public Library
% Source-File: IOI中国国家候选队论文/2006/陈启峰/陈启峰.ppt
\documentclass[11pt]{article}
\usepackage{oipl-vn}

\title{Thêm ràng buộc và nới điều kiện\\{\large Ghi chú theo slide}}
\author{Chen Qifeng}
\date{2006}

\begin{document}
\maketitle

\origtitle{Constraining and relaxing methods}
\sourcepath{IOI China candidate team papers / 2006 / Chen Qifeng / Chen Qifeng.ppt}

\section{Phạm vi}

Bộ trình chiếu đi cùng bài viết ``Một căng một chùng, con đường giải bài''.
Ghi chú này dựa trên bản bài viết và các lời giải phụ đã dịch.

\section{Hai thao tác tư duy}

Thêm ràng buộc là chủ động thu hẹp bài toán để cấu trúc trở nên rõ hơn, sau đó
chứng minh nghiệm trong miền hẹp vẫn đủ cho bài gốc. Nới điều kiện là tạm bỏ
bớt một yêu cầu khó, giải bài dễ hơn, rồi chứng minh có thể đưa nghiệm trở
lại bài ban đầu.

\section{Ba ví dụ}

\paragraph{Kỵ sĩ.}
Không gian bàn cờ vô hạn được thay bằng bài toán trên các vectơ sinh một lưới
nguyên. Ràng buộc mới giúp biến tìm kiếm hình học thành biến đổi số học.

\paragraph{Những loài vật thân thiện.}
Điều kiện thứ tự quá chặt được nới ra để bài toán có thể xử lý bằng quét và
cấu trúc dữ liệu đơn giản hơn. Sau đó chứng minh nghiệm thu được vẫn thỏa yêu
cầu ban đầu.

\paragraph{Trạm cứu hỏa.}
Bài này kết hợp cả hai hướng: nới điều kiện để đặt trạng thái quy hoạch động,
rồi thêm ràng buộc đủ mạnh để chuyển trạng thái có thể tính được.

\section{Ghi nhớ}

Khi điều kiện quá rối, hãy thử nới. Khi không gian quá rộng, hãy thử siết.
Điểm quyết định là chứng minh phép biến đổi không làm mất nghiệm cần tìm.

\end{document}
